%%
%% The abstract is a short summary of the work to be presented in the
%% article.
\begin{abstract}

Traversal-based graph analytics are a fundamental class of graph workloads involving systematically visiting vertices in a graph to explore its structure in a specific order. \textcolor{blue}{Breadth-First Search (BFS)}, \textcolor{blue}{Single-Source Shortest Path (SSSP)}, and \textcolor{blue}{Multiple-Source Shortest Path (MSSP)} are representative algorithms in this class. Traversal-based graph analytics are crucial in many fields including artificial intelligence, network analysis, social media, and many others.  We have conducted a comprehensive, performance- and measurement-driven study to understand the behavior, insights, and implications of traversal-based graph analytics on both multi-core \phantomsection\label{para:first-cpu-gpu}\textcolor{blue}{Central Processing Units (CPUs)} and \textcolor{blue}{Graphics Processing Units (GPUs)}. 
We introduce a set of performance profiling metrics, including instruction issue rate on both CPUs and GPUs; L3 cache miss rate and miss frequency on CPUs; and effective device memory bandwidth on GPUs, capturing most critical architectural factors to uncover program execution dynamics. 

In addition to BFS, SSSP, and MSSP, our study also includes the blossom algorithm, a foundational method to compute maximum matchings in general graphs.  X-Blossom, a recently proposed parallel framework of blossom algorithm, currently has an implementation only on multi-core CPUs. With our optimizations, its CPU performance improves by up to $3.26\times$. More importantly, we have developed a highly efficient GPU implementation of X-Blossom that achieves up to $39.09\times$ speedup over X-Blossom-Pro. 

While recent studies largely emphasize the massive parallelism of GPUs, not all applications can exploit it effectively.
Our intensive evaluation of four traversal-based graph primitives with distinct parallelism, locality, and computation patterns shows that the proposed profiling approach can characterize execution bottlenecks and determine architectural suitability between CPUs and GPUs.
We believe our profiling methodology is a practical approach for analyzing other parallel-processing applications to understand their performance with respect to the interaction between parallelism, data movement, and architectural resources.

% % Background 
% The blossom algorithm is a foundational method for computing maximum matchings in general graphs, supporting critical resource-allocation tasks in diverse domains such as logistics planning, large-scale clustering, and dynamic cloud scheduling. 
% As these workloads in modern applications scale to billion-edge graphs, it becomes essential to design high-performance algorithms to meet their stringent runtime and scalability requirements.
% % Issues with the SOTA work
% X-Blossom, a recently proposed parallel framework of blossom algorithm, makes an important step in this direction,
% but still leaves substantial performance headroom in its multi-core CPU implementation due to several inefficiencies.
% More fundamentally, CPUs are latency-oriented and therefore offer relatively few cores, which prevents X-Blossom’s irregular parallelism from being fully exploited.
% This naturally motivates offloading X-Blossom’s computation to modern throughput-oriented hardware accelerators such as GPUs for additional performance gains.
% % Summary
% In this paper, we present a comprehensive, performance- and measurement-driven study of X-Blossom that advances its performance frontier on GPUs and introduces a general methodology for evaluating algorithm–architecture compatibility and guiding CPU–GPU platform selection. 
% % Performance
% We first remove major inefficiencies in the existing multi-core implementation, producing an optimized CPU baseline, X-Blossom-Pro, which achieves up to 3.26× speedups. 
% Building on this variant, we develop X-Blossom++, a GPU-based implementation that fully exploits the framework’s inherent parallelism. Across a broad suite of large real-world graphs, X-Blossom++ consistently outperforms X-Blossom-Pro, achieving an additional speedup of up to 46×.
% % SIGMETRICS
% To understand the architectural suitability of X-Blossom, we use breadth-first search (BFS) as an offset workload and implement state-of-the-art BFS and X-Blossom on both CPUs and GPUs. 
% We then conduct extensive cross-platform measurements of all four implementations using two simple, architecture-agnostic metrics: instruction execution rate and effective memory bandwidth. 
% The contrast between BFS and X-Blossom reveals distinct metric signatures that explain why X-Blossom is best realized on GPUs. 
% From this analysis, we distill a general methodology that relates these metric characteristics to architectural suitability, offering practical guidance on when graph workloads should run on CPUs and when GPUs are the preferred platform.
    
\end{abstract}
